\section{Results}
\label{sec:results}

\subsection{Posterior summaries}
We illustrate a representative fit of the Echo Equation to the late-time dataset in Sec.~\ref{sec:data}. The finite-memory parameters remain small but non-zero, producing mild, phase-coherent modulations without degrading the overall goodness-of-fit relative to $\Lambda$CDM.

\begin{table}[t]
  \centering
  \caption{Illustrative posterior means $\pm$\,1$\sigma$ under the baseline dataset (see Sec.~\ref{sec:data}). Values are placeholders for layout; replace with your run outputs.}
  \label{tab:params}
  \begin{tabular}{lcc}
    \toprule
    Parameter & Mean $\pm$ 1$\sigma$ & Units \\
    \midrule
    $\tau$     & $0.16 \pm 0.06$ & Gyr \\
    $\kappa$   & $-0.015 \pm 0.005$ & -- \\
    $\omega_0$ & $1.9 \pm 0.3$ & rad per e-fold \\
    $Q$        & $14.0 \pm 4.5$ & -- \\
    $\Omega_m$ & $0.302 \pm 0.011$ & -- \\
    \bottomrule
  \end{tabular}
\end{table}

\subsection{Growth residuals}
A residual oscillatory structure appears in $f\sigma_8(z)$ with characteristic $\Delta z \simeq \pi/\omega_0$. Its phase aligns with the kernel frequency inferred from the background response, providing a cross-check across observables.
